% !PRINTED 490
\weberpage{490}

Eine algebraische Zahl gen\"ugt nicht nur einer, sondern unendlich vielen rationalen Gleichungen; denn multiplicirt man zwei beliebige Functionen von der Form $F(x)$ mit einander, so erh\"alt man eine Function derselben Form, die f\"ur $x = \Theta$ verschwindet, wenn einer der Factoren diese Eigenschaft hat.

Unter allen rationalen Gleichungen, denen eine algebraische Zahl gen\"ugt, ist eine von m\"oglichst niedrigem Grade, $f(\Theta) = 0$, wenn $f(x)$ die Form hat:
\weq{2}{f(x) = x^n + a_1 x^{n-1} + a_2 x^{n-2} + \cdots + a_n,}
und es kann auch nur eine solche Gleichung geben, wenn wir, wie bisher immer, den Co\"efficienten der h\"ochsten Potenz von $x$ gleich 1 annehmen.

Denn sind $f(x), f_1(x)$ zwei Functionen von der Form (2) von gleichem Grade $n$, so ist $f(x) - f_1(x)$ von niedrigerem als dem $n^{\text{ten}}$ Grade, und wenn sowohl $f(\Theta)$ als $f_1(\Theta)$ verschwindet, so verschwindet auch $f(\Theta) - f_1(\Theta)$; wenn also diese Differenz nicht identisch verschwindet, so gen\"ugt $\Theta$ einer Gleichung von niedrigerem als dem $n^{\text{ten}}$ Grade, was gegen die Voraussetzung ist.

\gesperrt{Die Function $f(x)$ ist im K\"orper der rationalen Zahlen irreducibel.}

Denn zerf\"allt $f(x)$ in zwei rationale Factoren $f_1(x)$ und $f_2(x)$, von denen jeder von niedrigerem Grade ist als $f(x)$, so gen\"ugt $\Theta$ einer der beiden Gleichungen $f_1(\Theta) = 0$, $f_2(\Theta) = 0$, was unserer Voraussetzung widerspricht.

\gesperrt{Ist $n$ der Grad der rationalen Gleichung niedrigsten Grades, der die Zahl $\Theta$ gen\"ugt, so nennen wir $\Theta$ eine algebraische Zahl $n^{\text{ten}}$ Grades.}

\begin{center}
\S.\,133.

\gesperrt{Ganze algebraische Zahlen.}
\end{center}

\gesperrt{Eine algebraische Zahl $\Theta$ wird eine ganze algebraische Zahl genannt, wenn sie einer rationalen Gleichung}
\weq{1}{\Theta^m + A_1 \Theta^{m-1} + \cdots + A_{m-1} \Theta + A_m = 0}
%CONT>>
